\chapter{Introduction}
Imagine a typical conference: you have a printed program with scheduled talks, audience members quietly sitting in rows, and experts with rehearsed presentations at the front. Now, instead imagine that you walk into a room with blank scheduling boards, no predetermined talks, and participants actively creating the conference agenda. This on-the-fly event is an Unconference.

These Unconferences have gained popularity in academic, technology, and professional communities for their ability to facilitate knowledge exchange. Unconferences have emerged as an alternative to traditional conference formats. Based on Open Space Technology principles, Unconferences follow a structure that involves collaborative agenda creation, breakout sessions, and collective sharing of insights. Chapter 2 goes further in depth on Unconferences.

However, Unconferences present challenges that existing event management tools designed for traditional conferences fail to address effectively. Although Unconferences rely on simple materials such as flipboards, markers, and post-it notes, there is a gap in digital solutions that can support the participant-driven and on-the-spot organizing of Unconferences. More specifically, scheduling sessions that capture participant preferences is a challenging problem, even for smaller Unconferences. Existing conference software does not seem to support this automated scheduling that captures participant preferences.

This thesis presents UnconfRS, a digital platform designed to support the needs of Unconference organizers and participants. Specifically, this work looks to answer whether modern scheduling techniques can be used to create an effective Unconference management tool and whether Rust is an effective language for implementing the backend and the scheduler. By addressing these questions, this work aims to provide a tool that can effectively manage Unconferences.